\section{Chương~\ref{sec:debugging}. Gỡ lỗi cỗ máy}

Điều điện cực nghe được, số chiều thấp của hoạt động, hiệu quả của các giao diện dùng mảng cứng, việc não và bộ giải mã học cùng nhau và các chấm sáng thị giác do kích thích đều đã được đo trong các thí nghiệm được bình duyệt; ước tính luồng dữ liệu là phép tính số học của chương; về thiết bị Neuralink cấy cho người, theo chỗ chúng tôi kiểm tra được, dữ liệu chủ yếu do chính công ty công bố.

{\footnotesize
\setlength{\tabcolsep}{3pt}\renewcommand{\arraystretch}{1.15}
\begin{longtable}{>{\raggedright}p{0.5cm}>{\raggedright}p{4.0cm}>{\raggedright}p{5.6cm}>{\raggedright}p{2.3cm}>{\raggedright\arraybackslash}p{1.7cm}}
\toprule
TT & Khẳng định & Thí nghiệm và điều đã đo & Nguồn & Trạng thái \\
\midrule\endfirsthead
\toprule
TT & Khẳng định & Thí nghiệm và điều đã đo & Nguồn & Trạng thái \\
\midrule\endhead
1 & Điện cực nghe xung của các nơ-ron trong bán kính cỡ một trăm micromét, thường từ một đến vài nơ-ron; chúng được phân biệt theo hình dạng & Chuột cống, vùng CA1, ghi đồng thời trong tế bào và ngoài tế bào: hình dạng xung ngoài tế bào lặp lại độ rộng và biên độ của xung trong tế bào. Ước tính: một tetrode có thể phân biệt $60$--$100$ nơ-ron, nhưng thực tế chỉ tách được khoảng sáu. & \cite{Henze2000extra} & đã đo \\
2 & Não có $8.6\cdot10^{10}$ nơ-ron; đầu dò nghe khoảng một phần trăm triệu & Đếm nhân tế bào trong huyền phù mô đã hòa tan: $86\cdot10^9$ nơ-ron trong não người. Tỉ lệ là phép tính số học của chương. & \cite{Azevedo2009} & đã đo \\
3 & Hoạt động của hàng trăm nơ-ron vỏ não vận động nằm gọn trong mười đến hai mươi biến chung & Tổng quan các bản ghi ở vỏ não vận động của khỉ: hoạt động của quần thể được mô tả bằng ít biến ẩn, chung cho nhiều nơ-ron. & \cite{Gallego2017} & đã đo \\
4 & Nhiễu của các nơ-ron lân cận tương quan, và một trăm nơ-ron mang gần bằng một nghìn (mệnh đề~\ref{prop:pooling}) & Vỏ não thị giác của khỉ: hệ số tương quan nhiễu của các nơ-ron lân cận khoảng $0.1$, và lợi ích từ việc lấy trung bình bão hòa ở một trăm nơ-ron. Lý thuyết về các tương quan giới hạn thông tin. & \cite{Zohary1994, MorenoBote2014} & đã đo \\
5 & Luồng thô $2\cdot10^8$~bit/s so với $8\cdot10^4$~bit/s trong xung~\eqref{eq:probedata} & Phép tính của chương theo dung lượng luồng xung~\eqref{eq:spikeentropy}. Thông số số hóa là thực: chip Neuralink là $19.3$~kHz, $10$~bit mỗi kênh. & suy ra trong chương; \cite{Rieke1997, Musk2019} & lý thuyết \\
6 & Hệ thống Neuralink đầu tiên: chip khuếch đại và số hóa, luồng thô đi qua cáp, xung do chương trình bên ngoài tách; tách trên chip, vỏ kín, vô tuyến và cấp năng lượng không dây là ở thiết bị cho người, theo công bố của công ty & Bài báo của công ty: toàn bộ luồng thô qua cáp USB-C, xung được một chương trình tách theo thời gian thực bên ngoài chip; nén trên thiết bị, cấp năng lượng và truyền không dây được nêu là kế hoạch cho thiết bị lâm sàng. Với thiết bị cấy cho người, chỉ có công bố của công ty. Việc tách xung trên chip là cần thiết là kết luận của chương từ~\eqref{eq:probedata}. & \cite{Musk2019}; báo cáo của Neuralink, không có bài báo bình duyệt & đã đo (bài báo của công ty); với người là báo cáo của công ty \\
7 & Tới $3072$ điện cực trên $96$ sợi, mỗi sợi $32$; các sợi được robot cắm, tránh mạch máu & Bài báo của công ty: $3072$ điện cực trên $96$ sợi, mỗi sợi $32$; robot cắm $6$ sợi ($192$ điện cực) mỗi phút; tới $70\,\%$ điện cực ở chuột cống được cấy mảng lâu dài nghe được xung. & \cite{Musk2019} & đã đo (bài báo của công ty) \\
8 & Ở người khoảng một nghìn kênh trên $64$ sợi; một phần sợi đã xê dịch; con trỏ khoảng $10$~bit/s & Chỉ có công bố của công ty. Tìm kiếm trong PubMed không thấy bài báo bình duyệt nào có kết quả trên người; điều này khớp với lưu ý trong chương. & báo cáo của Neuralink; không có bài báo bình duyệt & đã đo (báo cáo của công ty) \\
9 & Giao diện dùng mảng khoảng một trăm điện cực điều khiển được con trỏ & Người bị liệt, $96$ điện cực trong vỏ não vận động sơ cấp: ý định cử động tay làm thay đổi xung ba năm sau chấn thương; bộ giải mã cho ``con trỏ thần kinh'' (thư điện tử, tivi), điều khiển bàn tay giả và tay robot. & \cite{Hochberg2006} & đã đo \\
10 & Viết ``bằng bàn tay tưởng tượng'': khoảng $90$ ký tự mỗi phút, dưới $8$~bit/s & Người bị liệt bàn tay, cố gắng viết tay, bộ giải mã hồi quy: $90$ ký tự mỗi phút với độ chính xác $94.1\,\%$ theo thời gian thực, trên $99\,\%$ sau khi tự sửa lỗi. Ước tính theo bit là phép tính của chương. & \cite{Willett2021} & đã đo \\
11 & Những cố gắng nói được chuyển thành văn bản với tốc độ khoảng $60$ từ mỗi phút & Người đã mất khả năng nói rõ, các mảng điện cực trong vỏ não: $62$ từ mỗi phút; $9.1\,\%$ lỗi từ với bộ từ vựng $50$ từ, $23.8\,\%$ với bộ từ vựng $125\,000$ từ. & \cite{Willett2023} & đã đo \\
12 & Giao diện chỉ rút ra một phần nhỏ dung lượng: một nơ-ron có vài chục bit mỗi giây, một trăm nơ-ron có hàng nghìn (mệnh đề~\ref{prop:probe}) & Giới hạn trên~\eqref{eq:spikeentropy} là lý thuyết; so sánh với các dòng~8, 10 là kết luận của chương. Việc nút thắt là số chiều thấp và sự chưa hiểu mã, chứ không phải dây dẫn, chưa được kiểm tra bằng thí nghiệm trực tiếp. & \cite{Rieke1997} & lý thuyết \\
13 & Não và bộ giải mã học lẫn nhau & Khỉ điều khiển con trỏ; bộ giải mã tự điều chỉnh trong vòng kín, còn các nơ-ron đồng thời học lại: độ chính xác tăng, kỹ năng được giữ lại và ít bị ảnh hưởng bởi cử động của chính tay con vật. Lời khuyên tách tốc độ học là quy tắc kỹ thuật, trong chương không có thí nghiệm riêng. & \cite{Orsborn2014coad} & đã đo \\
14 & Dòng điện qua điện cực kích thích các nơ-ron và dây dẫn xung quanh bất kể loại & Vỏ não chuột nhắt và mèo, ghi canxi hai photon: ngay cả dòng yếu cũng kích thích các nơ-ron thưa thớt, rải rác trên khoảng cách tới vài milimét, thông qua kích thích trực tiếp các sợi trục trong một thể tích đường kính vài chục micromét. Tức là kích thích không chọn lọc, nhưng cũng không liền khối. & \cite{Histed2009stim} & đã đo \\
15 & Kích thích vỏ não thị giác làm bật các chấm sáng; tới $16$ chấm cùng lúc cho phép nhận ra một số chữ cái và đường biên đồ vật & Người mù, $96$ điện cực trong vỏ não thị giác trong $6$ tháng: ngưỡng của một điện cực $66.8\pm36.5$~\textmu A; kích thích đồng thời theo nhóm (tới $16$ điện cực, giới hạn của bộ kích thích) hạ ngưỡng và cho các mẫu phân biệt được; một số chữ cái và đường biên đồ vật được nhận ra. & \cite{Fernandez2021} & đã đo \\
16 & Hướng thứ hai của Neuralink là thiết bị đưa tín hiệu vào vỏ não thị giác & Chỉ có thông báo của công ty về việc phát triển; không tìm thấy dữ liệu bình duyệt về hoạt động của nó. & báo cáo của Neuralink & giả thuyết \\
17 & Nồng độ chất dẫn truyền thần kinh quảng bá có thể đo bằng cảm biến hóa học trong mô & Lát não chuột cống, đo volt-ampe vòng quét nhanh trên sợi carbon: đo được sự giải phóng và tái hấp thu serotonin; chất này thoát ra ngoài khớp thần kinh. & \cite{Bunin1998} & đã đo \\
\bottomrule
\end{longtable}}
